\section{The symmetric terminal levels and their Hermitian geometry}
\label{sec:terminal-geometry}

\status{TG1--TG12 construct the previously missing symmetric
penultimate level, terminal row chart, full and row auxiliary levels,
terminal incidence and Hermitian carriers. They finish the actual
\texttt{FormsTiedSym} and \texttt{ColsTiedSym} constructions.
They use the actual arithmetic
host and the stated \texttt{TC.Ok}, \texttt{Pins}, \texttt{Laws} and
\texttt{Later} data. Initial decorated-point realizers are supplied by
DR3, not retained as an additional hypothesis. Integral lattices,
specified residue maps, preferred class representatives and the common
global coordinate frame are preserved throughout.}

\paragraph{Conventions and supplied data.}
Keep the true column datum $I:\mathrm{TC}(B,X)$ and write
$D=I.\mathrm{rec7}(l,Z)$. The words $u$ range over the original word
set; all constructions below use the one auxiliary-prime datum $G$ of
FC13. Fix $u$ temporarily, put $p=T_u$, $\pi=\pi_u$ and
$\mathcal O=\mathcal O_E$, and use the standard coordinates of SF2:
\[
 h(x,y)=\bar x^{\mathsf T}y,\qquad
 \pi\bar\pi=p,\qquad(\pi,\bar\pi)=\mathcal O.
 \tag{TG.1}
\]
The selected type primes are $S=\{P:P\mid\pi\}$. Their residue
maps are the original maps $\operatorname{res}_P$; the selected
auxiliary primes lie over $g_0,g_1$, are disjoint from $S\cup\bar S$,
and satisfy $\pi\equiv1$ there. The reference twist has no $Q$-frame
primes. The supplied rank identities give
\[
 s=m/2-1,\quad r=k+2s+1=k+m-1,\quad
 N_{m-1}=r+(m-1)t,\quad n=N_m-1=r+mt.
 \tag{TG.2}
\]
Here $m\ge2$ is even and $t>0$ is odd. The row space is $E^n$,
the full terminal space is $E^{n+1}$, and the inclusion $\iota$ is
$x\mapsto(x,0)$. The distinguished vector spanning its orthogonal
line is not silently identified with the last basis vector.

\paragraph{TG1: the penultimate symmetric group is the actual residue group.}
For any relevant matrix size $d$, put
\[
 \rho(a)=
 \begin{cases}s+a&a<k,\\a-k&k\le a<k+s,\\a&a\ge k+s,
 \end{cases}
 \qquad
 \ell(a)=
 \begin{cases}\rho(a)&a<r,\\r+\lfloor(a-r)/t\rfloor&a\ge r.
 \end{cases}
 \tag{TG.3}
\]
The group $J_{\rm sym}(d)$ consists of the invertible residue
matrices $M$ such that both $M$ and $M^{-1}$ satisfy
\[
 M_{ab}=0\quad\hbox{if }\ell(b)<\ell(a),\qquad
 M_{aa}=1\ (a<r),\qquad M_{a,k+s}=0\ (a<k).
 \tag{TG.4}
\]
This is literally \texttt{JSymG} with its final Boolean equal to
\texttt{false}, hence the source's \texttt{HSym}. At step $m-1$
the actual finite group $\mathrm{levG}$, transported through the
specified auxiliary residue-field isomorphism, equals
$J_{\rm sym}(N_{m-1})$.

Here is a proof of both inclusions. At an auxiliary prime $P$ the
integer $q=g_a$ is a uniformizer in $\mathcal O_{E,P}$, since the
prime is completely split and has residue cardinality $g_a$.
The extra-space parameter at step $m-1$ is zero. Thus the actual
flag lattice at a head cut $b<r$, or a block cut $b=r+jt$, is
\[
 F_b=\{x:\ x_a\in q\mathcal O_{E,P}
                \text{ if }\rho(a)\ge d(b),\quad
                x_a\in\mathcal O_{E,P}\text{ otherwise}\},
 \quad d(b)=\begin{cases}s&b=k+s,\\b&b\ne k+s.\end{cases}
 \tag{TG.5}
\]
The marked vector for the head cut $b$ is $e_{\rho^{-1}(b)}$.
The local group preserves these lattices and satisfies
$(g-1)e_{\rho^{-1}(b)}\in F_b$. Taking $b=\rho(a)$ forces the
head diagonal to be one and every entry below it in the ordering
$\ell$ to vanish modulo $P$. At $b=k+s$ the cut begins at $s$;
its marked-vector condition additionally forces
$M_{a,k+s}=0$ for $a<k$. The block cuts force all lower off-block
entries to vanish. Applying the same argument to $g^{-1}$ proves
all of (TG.4) for the inverse as well: preservation of $F_b$ turns
$(g-1)e\in F_b$ into $(g^{-1}-1)e=-g^{-1}(g-1)e\in F_b$.

Conversely, lift $M$ entry by entry to $\mathcal O_{E,P}$.
Its determinant is a unit, so the lift $g$ and its inverse are
integral and reduce to $M,M^{-1}$. If a row lies on the scaled
side of a cut (TG.5) and a column does not, then
$\ell(\text{column})<\ell(\text{row})$. Therefore the corresponding
entry of $g$ is divisible by $q$. These divisibilities prove
$gF_b\subset F_b$, and the inverse conditions prove equality.
For a marked head vector, every scaled coordinate of
$(g-1)e_{\rho^{-1}(b)}$ vanishes modulo $P$: use the diagonal and
lower-entry conditions, and at the exceptional cut use precisely
the extra zero block in (TG.4). Division by the uniformizer now
gives the required membership in $F_b$. Hence the lift belongs
to the actual local symmetric flag group. The supplied
\texttt{grpG} and \texttt{last\_h} laws identify its reduction
image with $\mathrm{levG}$. This proves equality, including
surjectivity onto the stated finite group; it is not just a
containment of possible reductions.

\paragraph{TG2: the exact penultimate character and common level.}
For $b$ complete $t$-blocks after the head, define
\[
 \Lambda_b(M)=
 \psi\!\left(\sum_{\substack{0\le a<r-1\\a+1\ne k}}
                     M_{a,a+1}\right)
 \prod_{j=0}^{b-1}\chi_j\!\left(
       \det M_{[r+jt,r+(j+1)t)}\right).
 \tag{TG.6}
\]
This is exactly \texttt{lamSym}; the omitted edge is the entry
$(k-1,k)$ in the original coordinate order. The character of the
actual penultimate finite level is $\Lambda_{m-1}$, after the
same residue-field transport as in TG1.

Indeed, with extra-space parameter zero, the coordinate matrix
in the flag character is the original matrix. Its additive
argument is the two superdiagonal sums separated by the omitted
edge; its other factors are the displayed determinants. Reduction
commutes with these sums and determinants. The extra inverse
character has argument zero and contributes one. Every block
determinant is nonzero on the group, so the source's extension
of a multiplicative character by zero contributes its ordinary
unit value. FC13 identifies $\psi$ and every $\chi_j$ under the
prescribed residue embeddings. Agreement there on natural
integer residues suffices because each auxiliary residue field
is the prime field of order $g_a$. This proves (TG.6) on every
group element and also proves that it is a character: it is the
transport of the actual level character.

Take the common level $J$ and finite $J$-set of FC12, obtained
from the zero-coupling chart. TG1 identifies its auxiliary
reduction groups with $J_{\rm sym}(N_{m-1})$. FC14 then gives
the actual \texttt{GLevH}, with frame equal to the identity at
every auxiliary place and local character (TG.6), $b=m-1$.
The global character retains the source's inverse convention:
$\chi_J(x)$ is the product of $\Lambda_{m-1}$ evaluated at the
inverse reductions of $x$. None of the common-level construction
is changed. In particular
\[
 Y=\coprod_\beta J/U_\beta,\qquad
 \operatorname{Stab}_J([1]_\beta)=U_\beta,
 \qquad \chi_J|_{U_\beta}=\chi_\beta.
 \tag{TG.7}
\]
The finite reduction kernel lies in every $U_\beta$, which makes
all these coset sets finite. This proves the full
\texttt{stepBefore\_oneLevel} conclusion, including its specified
$\mathrm{HSym}$ and $\mathrm{lamSym}$ clauses.

\paragraph{TG3: construct the terminal row chart.}
The raw row action has empty type-prime set and, at the reference
twist, empty $Q$-frame set. Its valid decoration has a full
self-dual lattice $N_c\subset E^n$, the prescribed auxiliary
residue bases, and the actual finite groups
$\Lambda_P=\mathrm{rowTyG}_P$. Define the row base lattice to be
$\mathcal O^n$. At auxiliary places impose reduction in
$\Lambda_P$; elsewhere take the full integral unitary stabilizer.
Let $U_P$ be this adelic level and give it character
\[
 \chi_P(x)=\sigma\!\left(
       \prod_{Q\in S_G}\lambda_Q(\operatorname{red}_Q(x)^{-1})
                         \right).
 \tag{TG.8}
\]
Each preferred row decoration has an actual adelic representative
$y_P(c)$ by DR3, with these exact residue bases. Thus
$y_P(c)\widehat{\mathcal O}^{n}=\widehat N_c$.

The classes of the original row action identify bijectively with
the double quotient by the rational unitary group and $U_P$.
To check this assertion, write the group action as its rational
factor $\gamma$ and its auxiliary finite factor $a$. A relation
between two decorations gives, by the exact carried-lattice and
residue-basis equations, a matrix $x$ with
$\operatorname{rat}(\gamma)y_P(c')=y_P(c)x$ and $x\in U_P$.
Conversely this adelic equality recovers the lattice by rational
intersection (AD4), and recovers the finite factor by reduction;
it is the original decoration relation. For surjectivity, start
with any unitary finite adele, intersect its carried lattice
with $E^n$, and use its reduced inverse coordinates as the
auxiliary bases. At each place its lattice is a unitary image of
the self-dual standard lattice; the paired-dual calculation of
NB1 and the local-global intersection of LP1 make the global
lattice self-dual. Its reduced inverse coordinates are an actual
residue basis, so the resulting decoration is valid. Its preferred
representative is in the same double coset.

The stabilizer map is the actual rational factor $s\mapsto\gamma_s$.
It is injective: when $\gamma_s=1$, preservation of the auxiliary
bases forces every finite factor to be one; there are no other
factors because the reference $Q$-set is empty. Its range is
exactly the rational unitary elements in
$y_P(c)U_Py_P(c)^{-1}$. Equation (TG.8) is the corresponding
character relation, with the inverse already accounted for.
These statements prove every \texttt{ReadsP} and
\texttt{ClsPTied} field, not merely the existence of a name for
such a chart.

Apply CV10 to this actual chart, using its unchanged frame
identification. It gives the complex, rational and exact integral
row-value clauses. The constant shift of the terminal weight by
$m(t+1)/2$ leaves its polynomial weight, integral $\Delta$ lattice
and Fock pairing unchanged; the determinant factor stays in the
frame evaluation. FC7 therefore supplies the Petersson and
\texttt{NormPTie} clauses with the specified embeddings and
rational basis. We obtain the source's full
\texttt{clsP\_tie}, with its stipulated identity map on row
frame functions. No equality between integral $\Delta$ and
dual-Weyl lattices is used.

\paragraph{TG4: the full and row terminal levels and their frames.}
On $E^n$, the actual row finite group is the parabolic with
symmetric head and $m$ blocks of size $t$. By reading its
head and block conditions, it is $J_{\rm sym}(n)$ of (TG.4).
On $E^{n+1}$ the terminal finite group is the same parabolic
with one final block of size one, and with the final column
fixed. The lower block-zero condition forces the final row to
vanish away from its diagonal. Consequently it is exactly
$J_{\rm sym}^{+}(n+1)$, where both the final row and final
column are those of the identity, for the matrix and its inverse.
This is \texttt{JSymG true}. Conversely these conditions give
the original parabolic and fixed-column conditions, so both
identifications are equalities.

The character in each case is (TG.6) with $b=m$. The same
skipped-superdiagonal and block-determinant calculation proves
this, including the coefficient embedding. The full terminal
base lattice is
\[
 L_\pi=\mathcal O^n\oplus\pi\mathcal O,
 \tag{TG.9}
\]
whereas the row base is $\mathcal O^n$. At an auxiliary place
$v$, $\pi$ is a unit. For its prescribed split projection
$\phi_v$, choose
\[
 F_P(v)=I_n,\qquad
 F_M(v)=\operatorname{diag}(1,\ldots,1,\phi_v(\pi)).
 \tag{TG.10}
\]
The latter spans the same completed integral lattice as the
identity and commutes with the full residue group, because that
group fixes the last row and column. Thus the level description
and character are unchanged. FC14 supplies both actual
\texttt{GLevH} records in these frames. Their first columns
satisfy $F_M\iota=\iota F_P$, and their final column is
$\phi_v(\pi e_{n+1})$. These are exactly the two
\texttt{FrExtFr} equations. They hold for every auxiliary place
and use the one $G$ shared across all words.

\paragraph{TG5: the terminal vector, lattices and residue functional.}
Put $z_0=\pi e_{n+1}$. For the actual raw terminal extension of a
row decoration with self-dual lattice $N_c$, its three lattices
are distinct:
\[
 \widetilde N_c=\iota N_c\oplus\mathcal Oe_{n+1},\qquad
 \widetilde L_c=\iota N_c\oplus\pi\mathcal Oe_{n+1},\qquad
 \widetilde L_c^{\#}=\iota N_c\oplus\bar\pi^{-1}
                                      \mathcal Oe_{n+1}.
 \tag{TG.11}
\]
The raw type functional is
$\widetilde f_P(x)=\operatorname{res}_P(\bar\pi x_{n+1})$.
These identities also follow directly from the definitions:
the row type set is empty, so its $N,L,L^{\#}$ coincide;
the extra $N$ coordinate is integral; the simultaneous kernels
of $\operatorname{res}_P(\bar\pi x_{n+1})$ require
$x_{n+1}\in\bigcap_{P\in S}P=\pi\mathcal O$ by FC1;
and Hermitian duality sends $\pi\mathcal O$ to
$\bar\pi^{-1}\mathcal O$. Thus the normalization is
$\bar\pi$, rather than $\pi$.

The source's chosen movement $m_c$ takes the preferred full
representative of $\widetilde c$ to this raw extension. Let
$\gamma_c$ be its rational unitary factor. The finite factors
leave the type kernel lattice unchanged. Consequently
\[
 \gamma_c L_{\widetilde c}
      =\iota N_c+\mathcal O z_0,
 \qquad h(z_0,z_0)=\bar\pi\pi=p.
 \tag{TG.12}
\]
The vector $z_0$ spans the specified orthogonal line. Indeed,
any nonzero vector $z_c$ orthogonal to $\iota E^n$ has zero first
$n$ coordinates, so $z_0=(\pi/(z_c)_{n+1})z_c$. Nonvanishing
of the last coordinate follows from $z_c\ne0$. This proves the
required span statement without changing the distinguished
$z_c$. Applied to the standard row base, (TG.12) is also the
exact \texttt{TilBaseLat} identity.

For the functional clause, let $x\in L_{\widetilde c}^{\#}$ and
suppose $a=h(z_0,\gamma_cx)\in\mathcal O$. Set $b=\bar\pi$.
The dual-lattice identities at the selected and conjugate primes
give $bL_{\widetilde c}^{\#}\subset N_{\widetilde c}$.
On this latter lattice, FC5 identifies the quotient projection
with the original decoration projection, while the period
extension has exactly the functional in (TG.11). Hence
\[
 \operatorname{res}_P(b)\,
   f_P(\operatorname{pr}_P x)
   =f_P(\operatorname{pr}_P(bx))
   =\operatorname{res}_P(ba).
 \tag{TG.13}
\]
Here the rational movement is the same $m_c$ as in (TG.12);
its finite factors do not rescale the type functional.
Since $\operatorname{res}_P(b)\ne0$, cancellation gives
\[
 f_P(\operatorname{pr}_P x)
       =\operatorname{res}_P\bigl(h(z_0,\gamma_cx)\bigr).
 \tag{TG.14}
\]
This is the full \texttt{TilTied} functional clause on the
Hermitian dual, not just on its sublattice $N$. The identical
calculation on the base lattice, with FC6's exact completed
functional, gives \texttt{TilBaseFv}.

\paragraph{TG6: the actual terminal adelic incidence and evaluation.}
For a row unitary matrix $y$, let
$\operatorname{ext}(y)=\operatorname{diag}(y,1)$. By (TG.11)
and the integral adelic-span identity of AD4, this extension
carries $\widehat{\mathcal O}^{n+1}$ to
$\widehat{\widetilde N}_c$ whenever $y=y_P(c)$.
Let $k$ be the actual basepoint matrix of AD1. Then
$\operatorname{ext}(y_P(c))k$ realizes the raw terminal
decoration, including its exact residue bases. At a type prime
the inverse last row of $k$ has residue
$\bar\pi e_{n+1}$, which gives the functional in (TG.11).
At auxiliary primes $k=1$ and $\pi\equiv1$, so the new last
residue coordinate $x_{n+1}/\pi$ equals $x_{n+1}$ modulo the
prime; the other coordinates are the original row residue basis.
There are no reference $Q$-frames to check. This proves all
\texttt{IsPtOf} clauses of the extension.

Write $y_M(\widetilde c)$ for the chosen full chart representative.
The exact level-reading relation, applied to the finite part of
the same movement $m_c$ used in TG5, gives
\[
 x_c=y_M(\widetilde c)^{-1}
       \operatorname{rat}(\gamma_c)^{-1}
       \operatorname{ext}(y_P(c))\in U_M,
 \qquad \chi_M(x_c)=\sigma(\chi_{\rm col}(m_c)).
 \tag{TG.15}
\]
The terminal coupling set is trivial, so this reading introduces
no further chart condition. Set $\kappa_c=x_c^{-1}$. Then
\[
 \operatorname{rat}(\gamma_c)y_M(\widetilde c)
       =\operatorname{ext}(y_P(c))\kappa_c,
 \qquad
 \chi_M(\kappa_c)^{-1}=\sigma(\chi_{\rm col}(m_c)).
 \tag{TG.16}
\]
These equations fix both inverse conventions.

They prove the literal \texttt{TilAd} statement for every matrix
$g$ satisfying the stated extension conditions. In fact those
conditions force $g=\operatorname{ext}(y_P(c))$: the equation
$g\iota=\iota y_P(c)$ determines the first $n$ columns, and
$g z_0=z_0$ determines the last column because the image of
the nonzero field element $\pi$ is a unit in every coefficient
algebra under consideration. Thus no alternative extension is
being chosen in \texttt{TilAd}.

Finally, on frame functions the source's terminal restriction is
defined using $m_c$ as
\[
 (\operatorname{ev}_{z_0}\varphi)(c,z)
  =\sigma(\chi_{\rm col}(m_c))\,
           \varphi(\widetilde c,\gamma_c^{-1}z).
 \tag{TG.17}
\]
With the unchanged frame identification and $eW_z=\mathrm{id}$,
if $z=\gamma_c z'$ then (TG.16)--(TG.17) give
$(eW_z\operatorname{ev}_{z_0}\varphi)(c,z)
=\chi_M(\kappa_c)^{-1}\varphi(\widetilde c,z')$.
This proves the source's exact \texttt{EvzTie}, together with
the prescribed identity $eW_z=\mathrm{id}$. TG5--TG6 therefore
prove \texttt{til\_ties} with its actual chosen $\gamma_c$,
vector and lattices, not an abstract substitute incidence.

\paragraph{TG7: construct the full simultaneous symmetric forms.}
Choose $G$ once as in FC13. At every step and every word, DR3
supplies the decorated realizers used in FC3. Construct those
charts and apply FC4's functional normalization, which preserves
their lattice, local and global levels, characters, representatives
and stabilizer maps. Keep their original chart labels. The finite
class law and surjectivity of the label map make each label set
finite; the terminal label set is the one-element set, because
its coupling data are empty. At the endpoints reindex this
one-element chart as a \texttt{ClsTied}.

We now assign every part of \texttt{FormsTiedSym} to a construction
already proved. FC1 supplies the exact type-prime data. CV9 and
FC7 give the endpoint and middle complex, rational, integral,
Petersson and norm clauses. FC5--FC6 give their quotient maps
and normalized completed residue functionals; FC17 gives their
integral basis comparison. NB7--NB8, with DR3 discharging their
initial realizer premise, give all up and down naturality clauses.
FC18 gives endpoint and middle dilation, including preservation
of the chart label. FC14 gives the middle reduction groups and
characters, and FC15 gives the prescribed step-zero group.
TG1--TG2 give the common penultimate group, character, frame and
finite orbit carrier. TG3 gives the actual row class/value/norm
objects. TG4 gives the terminal full and row levels and their
frame extension. TG5--TG6 give every terminal lattice, functional,
adelic-incidence and evaluation clause. FC8 supplies the period
map and its metric compatibility with these same frame functions.

All frame identifications are the fixed coordinate identifications
of SF2; thus the required \texttt{DGam} equality and the identity
$eW_z$ are literal equalities. The quotient maps retain the
prescribed formulas on $N_c$ and on the down kernel lattice,
so \texttt{StdLvC} holds, including its identity penultimate
frame. The step-zero and row auxiliary frames are identity for
every word, and every $G$ is the one chosen above. Its
\texttt{CharacterCompatible} law is FC13. These assignments
construct the exact conclusion of
\texttt{formsTiedSym\_std}: a simultaneous family of full
\texttt{FormsTiedSym}, \texttt{StdF}, \texttt{StdLvC}, finite
chart sets, a unique terminal chart, a constant $G$, and the
stated character compatibility. No terminal field remains an
independent existence premise.

\paragraph{TG8: the discriminant, valuation orientation and projections.}
For every valid class lattice $L_c$, put $H_c=L_c^{\#}$.
Then
\[
                     H_c/L_c\simeq_{\mathcal O}\mathcal O/(p).
 \tag{TG.18}
\]
To prove this, use the self-dual companion $N_c$ and choose a
local basis at a selected prime $P$ adapted to its nonzero type
functional. Then $L_{c,P}$ is obtained from $N_{c,P}$ by
replacing one summand $\mathcal O_P$ by $P\mathcal O_P$.
At $\bar P$, $L_{c,\bar P}=N_{c,\bar P}$. Hermitian duality
between the paired primes therefore gives
$H_{c,P}=N_{c,P}$, and at $\bar P$ exactly one enlarged dual
summand. Consequently $(H_c/L_c)_P$ and
$(H_c/L_c)_{\bar P}$ are one-dimensional residue modules;
the quotient vanishes at every other prime. The same calculation
shows $pH_c\subset L_c$. FC1 gives
$(p)=\prod_{P\in S}P\bar P$ with no repetitions. The Chinese
remainder idempotents decompose the finite $\mathcal O/(p)$-module
$H_c/L_c$ into these one-dimensional factors. Choosing one
nonzero vector in each factor gives (TG.18). This proves the
source's \texttt{DiscCl}; it is a module equivalence, with no
claim of an isometry of discriminant pairings.

The source's \texttt{SeenCl} has the original valuation
orientation. For a coefficient place $w$, pull its valuation back
along the prescribed embedding
$E\xrightarrow{\sigma_{\tau(w,v)}}E\xrightarrow{j_E}L$;
call it $\nu$. The \texttt{tauOf} identity says that the positive
ideal of its restriction to $\mathcal O_F$ is the prescribed
$v$. The positive ideal on $\mathcal O_E$ is a nonzero prime:
all integral elements have nonnegative valuation, and $p$ has
positive valuation. It is therefore a maximal prime over $v$.
By FC1 there are exactly the paired primes $P,\bar P$ over $v$.
The supplied \texttt{seenP} identity gives $\nu(\pi)>0$ in
the seen case, hence selects $P$; in the other case it gives
$\nu(\bar\pi)>0$ and selects $\bar P$. Thus
\[
 \nu(a)>0\quad\Longleftrightarrow\quad
 \begin{cases}a\in P&\text{if seen},\\a\in\bar P&\text{otherwise}
 \end{cases}\qquad(a\in\mathcal O).
 \tag{TG.19}
\]
Injectivity of the distinguished coefficient embedding identifies
the $L$-element appearing in \texttt{SeenCl} with this precise
pullback. This rules out changing either the residue prime or
the embedding in (TG.19).

Let $C_i$ denote the fixed coordinate isometry from the original
Hermitian space to $E^{N_i}$. Directly from the pairing,
$C_i(L^{\#})=(C_iL)^{\#}$. Linearity also gives
$C_i(P L)=P C_i(L)$, commutation with sums, and commutation with
localization: a witness $az\in L$, $a\notin P$, is sent to the
same witness $aC_i z\in C_iL$. Therefore the Hermitian
projections may be defined by
\[
 \operatorname{pr}^{H}_P(z)=\operatorname{pr}_P(C_i z),
 \qquad
 \operatorname{prk}^{H}_P(z)=\operatorname{prk}_P(C_i z).
 \tag{TG.20}
\]
They have exactly the surjectivity, residue-linearity and kernels
of FC5. In particular the two kernels are $PH_c$ and
$P(L_c+PH_c)$, respectively; the kernel of the composed type
functional is $L_c+PH_c$. The multiplication-by-$p$ relation
between the two quotients is preserved by $C_i$.

At a middle step the standalone and chart projections are equal
on the whole required domains. They agree by definition on
$N_c$ and on the original down lattice. For arbitrary
$z\in H_c$, multiply by $b=\bar\pi$; then $bz\in N_c$.
For $z\in L_c+PH_c$, the same multiplication lies in the
original down lattice, by its local kernel description at $P$
and the unchanged conditions at other primes. Apply the known
projection equality there and cancel the nonzero scalar
$\operatorname{res}_P(b)$. This proves both equalities in
\texttt{pr\_lvC\_eq\_lv}, including their integral domains.

\paragraph{TG9: Hermitian row carriers and every level-prime move.}
Pull each row lattice of TG3 back by $C_P=\operatorname{crdP}$,
and transport its unitary level and rational stabilizer matrices
through the same rank reindexing. The row base becomes the
original standard lattice $X'.L_0$, since its orthonormal frame
has integral coordinates. This gives the row class carrier,
with identical representatives, double quotient and stabilizer
range. Its level is hyperspecial away from the two auxiliary
primes. It contains the required principal congruence subgroup
at an auxiliary place: if $(x-1)$ sends the standard integral
lattice into $g_a$ times itself, testing on coordinate vectors
gives a matrix $x=I+g_a Z$ with $Z$ integral. Its residue is
$I$, which lies in the row residue group of TG4. Thus the
literal \texttt{StepLevel0}, \texttt{ClassTie} and
\texttt{SameCar} conditions all hold. The lattice identity is
$C_P(\mathrm{lattice}_2(c))=N_c$.

Transport the actual NB up and down matrices by
$\Gamma=C_i^{-1}\gamma C_i$. Write $L=L_c$, $H=L^{\#}$,
$K=L+PH$, and let primes and quotient vectors keep their
original meanings. For an up vector $z\in H$ with
$f_P(\operatorname{pr}_Pz)=1$, NB5 gives
\[
 (\gamma L_{c'})_P=H_P,\qquad
 (\gamma H_{c'})_P=(H+\mathcal O p^{-1}z)_P,
 \quad f'_{P}(\operatorname{pr}'_P x)
      =f_P(\operatorname{pr}_P(p\gamma x)).
 \tag{TG.21}
\]
The last identity is asserted on the exact domain
$x\in H_{c'}$, $p\gamma x\in H$. For a down functional
$g\ne0$, the corresponding formulas are
\[
 (\gamma H_{c'})_P=L_P,\qquad
 f'_{P}(\operatorname{pr}'_P x)
             =g(\operatorname{prk}_P(\gamma x))
       \quad(x\in H_{c'},\ \gamma x\in K).
 \tag{TG.22}
\]
NB6 gives in both cases the same moving matrix and a factorization
\[
       \operatorname{rat}(\gamma)y_{c'}=y_c\eta\kappa,
       \quad\kappa\in U_\beta,
       \quad\eta_v=1\text{ for }v\ne P\cap\mathcal O_F.
 \tag{TG.23}
\]
DR3 supplies its realizer hypothesis. At any other place,
(TG.23) and the level condition show that $\gamma$ carries the
new completed lattice to the old one and preserves its
normalized functional. Rational intersection then gives equality
of localizations off $P,\bar P$; FC1 identifies these as the
only two primes above the support place. The other type
functionals agree by the same completed equation and FC6.
Using (TG.20), and cancelling the coordinate isometry, proves
all clauses of the literal Hermitian \texttt{UpCl} and
\texttt{DnCl}.

The diamond case is equally constructive. For $a\in k_P^\times$,
choose an integer representing $a$. At the split base place,
take a unit whose residue is $a^{-1}$ and multiply the last
coordinate in the selected split factor by that unit and in
the other factor by its inverse. Set the adelic matrix $\delta$
to be the identity elsewhere. It is unitary, preserves the base
lattice, and satisfies $a f(\delta x)=f(x)$; thus it is an
actual \texttt{IsDia} element. FC18 gives the exact diamond
class and a factorization (TG.23). Both $\delta$ and $\kappa$
stabilize the completed base lattice, so AD4 gives
\[
 \gamma L_{c'}=L_c,\qquad
 f'_{P}(\operatorname{pr}'_P x)
        =a f_P(\operatorname{pr}_P(\gamma x)),
 \tag{TG.24}
\]
with equality of every other type functional. These are the
\texttt{DilCl} clauses after coordinate transport.

NB4 and FC18 preserve the middle chart label. At the terminal
step its uniqueness supplies the same assertion. Equation
(TG.23), reindexed without changing its factors, is consequently
the required \texttt{Ad0S} or \texttt{AdIS} relation for each
move. This proves all six clauses of
\texttt{SymmetricAdelicCompatibility}. For indices beyond $m$,
or move parameters not satisfying their defining nonzero
condition, choose the identity isometry; the source asserts no
move identity in those cases.

\paragraph{TG10: the terminal Hermitian carrier and stabilizer map.}
Define the Hermitian type-prime datum $A_u$ using the original
$\pi$, primes, residue maps and coefficient-place orientation
from FC1 and TG8; use projections (TG.20), the move isometries
of TG9, and the row carrier just constructed. Put
\[
 z_0^H=C_m^{-1}z_0,\qquad
 \widetilde\Gamma_c=C_m^{-1}\gamma_c C_m.
 \tag{TG.25}
\]
Equations (TG.12)--(TG.14), together with
$C_m\iota_H=\iota C_P$, prove the Hermitian terminal lattice
and functional clauses. Equation (TG.18) gives its discriminant
clause and (TG.19) its seen-prime clause. Coordinate transport
proves all remaining quotient and multiplication clauses as in
TG8. Thus $A_u$ is the actual \texttt{AtT}, not just selected
maps without their laws.

Every one of the nine \texttt{CarT} requirements is now fixed:
$\pi$ and the residue map are unchanged; the row carrier and
row-lattice identity are TG9; the two endpoint projections are
(TG.20); the middle projection is TG8's cancellation identity;
and the vector and matrix are (TG.25). This also proves the
source's \texttt{hermitianSide\_classes\_eq\_formClasses}
constructor. The same coordinate construction applies to any
supplied $H$ with the literal \texttt{CarH} equalities, so the
\texttt{hermitian\_data\_at\_level\_prime} conclusion does
not require an additional choice of compatible carriers.

For the stabilizer clause, let $s$ stabilize the preferred row
class $c$. Its raw extension has rational matrix
$\operatorname{diag}(\gamma_s,1)$, so it intertwines $\iota$
and fixes $z_0$. The source's stabilizer map is conjugation by
the actual movement $m_c$:
\[
 \operatorname{stabMap}_c(s)=m_c^{-1}\operatorname{ext}(s)m_c.
 \tag{TG.26}
\]
Taking rational factors yields
\[
 \gamma_c\,\gamma_{\operatorname{stabMap}_c(s)}\,
                      \gamma_c^{-1}
       =\operatorname{diag}(\gamma_s,1).
 \tag{TG.27}
\]
Transport (TG.27) through $C_m,C_P$. On the hyperplane it is
the row stabilizer action, and on $z_0^H$ it is the identity.
These are exactly the two
\texttt{SymmetricStabilizerMapCompatibility} equations.

For the source's formulation with an arbitrary supplied
\texttt{FormsTiedSym} satisfying \texttt{StdF}, the same
$\gamma_c$ is forced, rather than assumed. Its evaluation
clause has $eW_z=\mathrm{id}$ and the standard frame map.
Test that clause on the indicator of one invertible frame.
The nonzero character scalar forces
$\operatorname{frMul}(\gamma_{m_c}^{-1}\gamma_c)z=z$.
Cancel the invertible frame and use injectivity of one prescribed
field embedding to obtain $\gamma_c=\gamma_{m_c}$. Thus
(TG.27) holds in that formulation too. With TG9, this proves
the full source conclusion: an $A$ with \texttt{CarT},
symmetric adelic compatibility and symmetric stabilizer
compatibility, under the stated supplied data.

\paragraph{TG11: normalize the level-zero functional and all residue ties.}
There is one functional choice in passing from forms to a
Hermitian level carrier that \texttt{SameCar} alone does not
specify. We make it explicit. At level zero replace the
Hermitian target module and functional by the corresponding
forms target module and functional, pulled back through the
fixed rank reindexing. This preserves its base lattice, every
local and global level, representatives, stabilizer maps,
auxiliary splittings and auxiliary frames. The stabilizer of
the functional is preserved because the forms level has
exactly that stabilizer, and \texttt{SameCar} identifies its
local group. Nonvanishing, cardinality, the $p$-annihilation
law and vanishing on the base lattice are transported as well.
Thus this is a valid replacement of a \texttt{ColLevel},
with all the equalities required by
\texttt{stepZeroHermitianLevel\_with\_formsFunctional}.

The new target module $M_v$ is annihilated by the actual
local image of $\pi$. To check it without assuming an extra
linearity property of the additive equivalence $M_v\simeq k_P$,
choose a class and a dual-lattice vector whose projected type
functional is nonzero. Such a vector exists by surjectivity of
the projection and nonvanishing of that functional. Let
$m_0\ne0$ be the corresponding value of the completed base
functional. Multiplying the vector by $\pi$ multiplies its
projection by $\operatorname{res}_P(\pi)=0$, so the same
functional comparison gives $\pi m_0=0$. The additive group
of $M_v$ has prime order $p$; every nonzero element generates
it. Since scalar multiplication commutes with integer
multiplication, $\pi M_v=0$. This is precisely the source's
\texttt{pi\_kills}, and proves the replacement's additional
annihilation conclusion.

For clarity, the residue comparison can also be constructed for
any carrier satisfying that annihilation conclusion, not only
the replacement just chosen. At $v\mid p$ write
$S_v=\operatorname{intV}(E,v)$. The prescribed split orientation
identifies it with the product of the two complete local integer
rings. The selected component of $\pi$ is a uniformizer and
the other is a unit. Consequently
\[
 S_v/(\pi)\simeq\mathbb F_p,\qquad
 L_{\pi,v}^{\#}/L_{\pi,v}\simeq S_v/(p),
 \tag{TG.28}
\]
where the second quotient is generated over $S_v$ by the
class of $\bar\pi^{-1}e_{\mathrm{last}}$. Both nonzero completed
functionals in question vanish on $L_{\pi,v}$, take values
in groups of order $p$, and are killed by $\pi$ on the target.
They therefore factor through the one-dimensional
$S_v/(\pi)$-quotient of (TG.28). Each induced map is an
isomorphism of additive groups. Composing one with the inverse
of the other gives the unique additive equivalence of their
targets which identifies the two functionals. This proves the
source's cyclic-kernel argument directly, with its selected
prime orientation; an arbitrary abstract group of order $p$
would not determine that orientation by itself.

The resulting $fMe_v$ gives the exact level-zero residue identity
on every class: if $y_c t=C_0z$ locally, then
\[
        f_P(\operatorname{pr}^{H}_P z)=fMe_v(fT_v(t)).
 \tag{TG.29}
\]
For intermediate charts define the residue functional on their
completed Hermitian dual by first applying the coordinate
reindexing, then the forms functional, then its prescribed
$fMe$. The equation $y_ct=C_i z$ is equivalent, after this
reindexing, to $t=y_c^{-1}C_i z$. FC6 and TG8 therefore prove
the identical relation (TG.29) there. At the terminal step the
unique chart identifies this with its endpoint functional.
These are all the \texttt{SymmetricResidueFieldLevels}
clauses; no linear-algebra existence theorem for an unspecified
functional is being substituted for their pointwise equations.

Finally set the level-zero auxiliary generators to the original
$X.\mathrm{piG}$. Their designated-subfield membership, norms,
split orientation and scalar-diagonal splitting laws are FC13.
The frame-lattice equivalence is FC15 transported by the same
rank reindexing, and the last two fields are the annihilation
law and (TG.29). This constructs the literal \texttt{ColLevelX}
for every word, hence \texttt{SymmetricLevelZeroData}, with
the exact identity $\mathrm{piG}=X.\mathrm{piG}$ for every word.

\paragraph{TG12: the complete symmetric column comparison.}
Construct the forms family by TG7. FC9 transports its charts to
a Hermitian carrier $H$, and TG8--TG10 construct $A$ with
\texttt{CarT} and both symmetric compatibility laws. Now apply
TG11's level-zero functional replacement. Every \texttt{CarH}
field survives, as do $\mathrm{CarT}$ and the two compatibility
laws: their level-zero clauses use only the unchanged group,
representatives and stabilizer maps. The type-prime indexing,
residue fields, other levels and class data are also unchanged,
so the same $A$ remains well typed. This order makes clear that
the functional normalization is not an input to its own
Hermitian carrier construction.

Two further displayed identities finish the Hermitian column
record. The supplied \texttt{Pins.cof} equation, with the
unchanged $\pi_u$, is exactly
\[
                 c_u\pi_u=\pi_{u_0}.
 \tag{TG.30}
\]
For each endpoint or chart base lattice, SF2 and the standard
base-lattice formula give $L_\pi$ as in (TG.9), in the
appropriate dimension. Its dual is
$\mathcal O^{d-1}\oplus\bar\pi^{-1}\mathcal O$, so
coordinatewise
\[
               L_\pi+\bar\pi L_\pi^{\#}=\mathcal O^d.
 \tag{TG.31}
\]
This is SF3's companion identity, and the right side is the
coordinate image of the original standard lattice. Equation
(TG.31) is precisely \texttt{SymmetricBaseCompatibility},
including the version with the principal ideal
$(\bar\pi)$ multiplying the dual. Together with TG11, these
prove the actual remaining conclusions of
\texttt{herm\_rest\_tied}: cofactor compatibility, base
compatibility, all residue-field levels and level-zero data
with the prescribed auxiliary generators.

Take these objects as the fields of \texttt{ColsHermS}, with
SF2's \texttt{StdTied}, and use the TG7 forms as its forms
family. The carrier equalities are FC9 and TG10; all projection,
vector and rational-matrix equalities are (TG.20), (TG.25) and
their stated chart versions. The splitting equality and
step-zero frame equality come from the same FC9 transport.
All auxiliary generators agree because TG11 used
$X.\mathrm{piG}$. All $G$ objects agree because FC13 was chosen
once. Their character compatibility is FC13. The step-zero
frames are identity by FC15, and the row frames are identity
by TG4, so both required frame-constancy assertions hold for
every pair of words.

This constructs the source's exact endpoint
\[
 \exists\,C:\mathrm{ColsTiedSym}(D),\qquad
       I.\mathrm{StdC}(l,Z,C)\ \land\
       \forall u\ I.\mathrm{StdLvC}(l,Z,u,C.\mathrm{forms}(u)).
 \tag{TG.32}
\]
It is the conclusion of \texttt{colsTiedSym\_std}. Each
substantive supplier in this final join has been given above
or proved in the cited earlier claims. In particular neither
\texttt{FormsTiedSym}, \texttt{CarT}, initial decorated-point
realizability, nor terminal incidence remains an independent
input to (TG.32). The conclusion still concerns the supplied
arithmetic host and true-column data with their stated laws
and \texttt{Later}. It does not construct that arithmetic
host or the later analytic package, and does not by itself
assert RH or GRH. The stronger \texttt{OneLevelSymmetricColumns}
assertions about all earlier chart groups and singleton chart
sets are separate from (TG.32).

